% =====================================================================
%  Chapter 19 opener -- "Chapter at a Glance", three reading paths + the
%  ladder staircase, THE running examples (the 16-slot array A shared with
%  Ch.18, this time folded into a segment tree of intervals -- the sum,
%  min and max trees over the SAME leaves -- plus the 10-node tree T of
%  Ch.12-17 with its Ch.16 preorder tin for the 19.5 Euler pointer), the
%  "which aggregate / which combine" cost chart, the who-owns-what
%  comparison (Ch.5 owns flattening; Ch.18 owns the BIT; Ch.20 owns lazy;
%  Ch.21 owns merge-sort/persistent/sparse-table; Ch.22 owns HLD; THIS
%  chapter owns the interval-binary-tree range engine), recognition guide,
%  constraint->tool table, notation.
%  Exports fig:c13-tree and fig:c13-taxonomy.
%  Hand-computed values the whole chapter reuses (verified once here).
%  Running array (16 slots, first 8 = the CSES samples, shared with Ch.18):
%    A = (3,2,4,5,1,1,5,3,7,0,2,4,1,6,3,2),  sum A = 49.
%  Recursive heap layout: root t[1], children t[2v],t[2v+1]; because n=16
%  is a power of two this complete tree has 31 nodes and its leaves are
%  t[16..31] = A[1..16].  The three folds over the internal nodes t[1..15]:
%    node:  t1   t2   t3   t4  t5  t6  t7  t8 t9 t10 t11 t12 t13 t14 t15
%    intvl: 1-16 1-8  9-16 1-4 5-8 9-12 13-16 1-2 3-4 5-6 7-8 9-10 11-12 13-14 15-16
%    SUM:   49   24   25   14  10  13  12   5  9   2   8   7   6   7   5
%    MIN:   0    1    0    2   1   0   1    2  4   1   3   0   2   1   2
%    MAX:   7    5    7    5   5   7   6    3  5   1   5   7   4   6   3
%  The iterative 2n bottom-up form (19.4) uses the SAME array here: leaves
%  t[16..31]=A, internal t[i]=combine(t[2i],t[2i+1]) -- identical values
%  because n is a power of two; the two forms differ only in the query walk.
%  Shared tree T of Ch.12-17 (edges 1-2,1-3,2-4,2-5,3-6,3-7,5-8,7-9,9-10),
%  Ch.16 preorder tin=(1,2,6,3,4,7,8,5,9,10), INCLUSIVE tout=(10,5,10,3,5,
%  7,10,5,10,10): subtree(2)=[2,5]={2,4,5,8}, subtree(7)=[8,10]={7,9,10}.
%  Every overlay below is hand-computed; every later section re-uses exactly
%  these numbers -- if a section disagrees with this opener, the OPENER is
%  authoritative: fix the section, never the other way.
% =====================================================================

\begin{chaptertheorybox}{Chapter at a Glance}
Chapter~18 ended with a boundary it could not cross: a Fenwick tree answers
\emph{sums and counts}, but the moment a range operation stops being
invertible -- a minimum under arbitrary point \emph{writes}, a greatest common
divisor, a maximum after decrements -- the BIT has no answer and hands the
range to \chref{c13}.  This chapter builds that structure.  A \term{segment
tree} is a binary tree whose nodes are \emph{intervals} of the array: the root
owns \([1,n]\), each internal node splits its interval at the midpoint, the
leaves own single positions, and every node stores one \emph{aggregate} of its
whole interval.  A single \term{merge} rule (associative, with an
\term{identity element}) turns \code{build}, \code{update} and \code{query}
into three \(\bigO{\log n}\) recursions -- and, unlike the BIT, the merge can
be almost any fold you like, including one that keeps a four-field tuple.

\textbf{You will learn to}
\begin{itemize}[nosep]
  \item see the interval hierarchy, write the \code{build}/\code{update}/
        \code{query} recursions, and choose a merge plus its identity element;
        \hfill\lv{2}--\lv{3}
  \item answer point-set + range-min end to end (CSES~1649, \(\secref{c13-core}\))
        -- the write-heavy query a BIT structurally cannot serve; \hfill\lv{3}
  \item carry a \emph{tuple} in each node and merge non-trivial aggregates
        (sum/prefix/suffix/best) in \(O(1)\) for the maximum-subarray combine
        (worked on CSES~1190, \(\secref{c13-combine}\)); \hfill\lv{4}
  \item \emph{descend} the tree to find the first position past a threshold,
        contrasting the walk with Ch.18's lowbit-lift (CF~1179C,
        \(\secref{c13-search}\)); \hfill\lv{4}--\lv{5}
  \item write the compact iterative (bottom-up) \(2n\) tree and read exactly
        when it beats the recursive form (\(\secref{c13-iter}\), CF~514D); \hfill\lv{4}
  \item recognise the static/offline range problems the tree nails, and know
        which of them Ch.20's lazy tags or Ch.21's sparse table own outright
        (\(\secref{c13-static}\)). \hfill\lv{5}
\end{itemize}

\textbf{Prerequisites.} Chapter~18 (\chref{c14}) for the \emph{contrast}: this
chapter repeatedly asks ``could a BIT do this?'' and answers it, so having the
BIT's prefix/lift machinery in hand sharpens every boundary.  Chapter~1
(\chref{c01}) for binary search as the idea a tree-descent re-does on the
aggregate.  Chapter~5 (\chref{c24}) only for the ``subtree \(=\) interval''
reduction the \(\secref{c13-static}\) Euler pointer consumes by label.  Recursion
and the array-as-heap layout (a node \(v\) with children \(2v,2v{+}1\)) are all
that is assumed.
\end{chaptertheorybox}

% ---------------- Reading paths + ladder position ----------------
\begin{chaptertheorybox}{How to Read This Chapter (three paths)}
Level badges \(\lv2\)--\(\lv5\): \(\lv2\) is reading intervals off a drawn
tree and folding a leaf row up to the root by hand; \(\lv3\) is a working
point-update / range-query tree on a real input (CSES~1649 is exactly this);
\(\lv4\) is choosing \emph{what a node stores} -- a scalar, a tuple, a count --
\emph{before} writing the merge, and picking recursive vs iterative; \(\lv5\)
is descending the tree for a threshold answer, or recognising that a
``static RMQ'' is a job for Ch.21's sparse table while a
\emph{mutable} one is a job for this chapter.

\textbf{Ground-up path (\(\lv2\)--\(\lv4\), ``I have never built a segment
tree'').} \secref{c13-core} in full: the interval picture, the three
recursions, the identity element, CSES~1649 with every build/update/query
step traced on \(A\).  Then \secref{c13-combine}: the tuple-node idea, whose
sharpest example (maximum subarray, CSES~1190) is a query no BIT can even
state.

\textbf{Query-hammer path (\(\lv4\)--\(\lv5\), ``the answer is a threshold
search'').} \secref{c13-core} for the mechanics only, then
\secref{c13-search}: run the tree \emph{downward} instead of summing it, so
``first element \(>x\)'', ``first prefix \(\ge\mathit{target}\)'' and ``k-th
by counts'' each cost \(\log n\).  Read it beside Ch.18's \(\secref{c14-order}\)
lifting -- the two are the same search seen from opposite ends of the tree.

\textbf{Compact / static path (\(\lv3\)--\(\lv5\), ``smaller code, faster
constant'').} \secref{c13-iter}: the \(2n\) bottom-up tree, its half-open
\([l,r)\) query, and the honest BIT-vs-segtree table that tells you which
structure to reach for.  Finish at \secref{c13-static} -- the offline
``first \(\ge x\)'', the RMQ boundary (where a sparse table wins and you
should stop here and go to Ch.21), and the light Euler-on-a-tree pointer that
Ch.5 opened and Ch.20/22 will finish.
\end{chaptertheorybox}

\begin{figure}[H]
\centering
\begin{tikzpicture}[diagram, x=9mm, y=7mm]
  \foreach \st/\lab in {1/lv1,2/lv2,3/lv3,4/lv4,5/lv5,6/lv6} {
    \fill[\lab!22, draw=\lab, line width=0.8pt] (\st-1,0) rectangle (\st,\st*0.52);
    \node[font=\scriptsize\bfseries, text=\lab] at (\st-0.5,\st*0.52-0.42) {L\st};}
  \node[annot, anchor=south west, font=\scriptsize]
    at (0.05,0.62) {19.1 the interval tree; fold leaves to root; merge + identity};
  \node[annot, anchor=south west, font=\scriptsize]
    at (1.05,1.14) {19.1 build/update/query recursions; solved CSES 1649};
  \node[annot, anchor=south west, font=\scriptsize]
    at (2.05,1.66) {19.2 tuple nodes: max-subarray merge (CSES 1190), gcd, count};
  \node[annot, anchor=south west, font=\scriptsize]
    at (3.05,2.18) {19.3 descend the tree: first \(>\)x, first prefix \(\ge\)t (CF 1179C)};
  \node[annot, anchor=south west, font=\scriptsize, text=lv5]
    at (4.05,2.70) {19.4 iterative 2n form + the honest BIT-vs-segtree table (CF 514D)};
  \node[annot, anchor=south west, font=\scriptsize, text=lv6]
    at (5.05,3.22) {19.5 static/offline ranges, RMQ boundary, Euler pointer + Summit};
\end{tikzpicture}
\caption*{Chapter~19 on the six-step ladder: reading and folding the tree is
\(\lv2\)--\(\lv3\); choosing what a node stores (scalar, tuple, count) and the
iterative form is \(\lv4\); descending for a threshold and honouring the
static/dynamic boundary is \(\lv5\); composing with offline sweeps and the
Summit is step~6.  Ground-up path \(=\) steps 2--4; full contest path runs to~6.}
\end{figure}

% ---------------- The running examples ----------------
\subsection*{The Running Example: One Array, Three Folds}

\emph{One} array carries the whole chapter -- the same 16-slot \(A\) as
Chapter~18, whose opening half is the sample array of CSES's range-query tasks
(1648/1649/1650 all print \(3,2,4,5,1,1,5,3\)), extended to a power of two so
the tree is complete:
\[
  A=(\,3,2,4,5,1,1,5,3,\;7,0,2,4,1,6,3,2\,),\qquad \textstyle\sum A = 49 .
\]
Chapter~18 stored \(A\) in a Fenwick array of \emph{lowbit blocks}; this
chapter stores the very same \(A\) in a binary tree of \emph{intervals}.  The
leaves are the array, and each node above is a fold of everything under it.  We
fold \(A\) three ways -- sum, min, max -- so that \figref{c13-tree}'s single
picture and the table below fix every number the sections will reuse:
\[
  \text{root: }\ \textstyle\sum=49,\quad \min=0,\quad \max=7 .
\]

\begin{figure}[H]
\centering
\begin{tikzpicture}[diagram, scale=0.46, transform shape]
  % ---- leaf row A[1..16] ----
  \foreach \i/\vv in {1/3,2/2,3/4,4/5,5/1,6/1,7/5,8/3,9/7,10/0,11/2,12/4,13/1,14/6,15/3,16/2} {
    \draw[primaryblue, line width=0.6pt, fill=primaryblue!6] (\i-0.5,-0.3) rectangle (\i+0.5,0.5);
    \node[font=\scriptsize\ttfamily] at (\i,0.1) {\vv};
    \node[font=\tiny\ttfamily, text=darkgray] at (\i,-0.55) {\i};}
  % ---- internal nodes t8..t15: box labelled t_i=value, interval below ----
  \foreach \n/\x/\s/\iv in {
      8/1.5/5/{1\text{--}2}, 9/3.5/9/{3\text{--}4}, 10/5.5/2/{5\text{--}6}, 11/7.5/8/{7\text{--}8},
      12/9.5/7/{9\text{--}10}, 13/11.5/6/{11\text{--}12}, 14/13.5/7/{13\text{--}14}, 15/15.5/5/{15\text{--}16}} {
    \node[boxnode, primaryblue, minimum size=8mm, font=\scriptsize] (t\n) at (\x,1.4) {\(t_{\n}{=}\s\)};
    \node[font=\tiny\ttfamily, text=green!45!black, anchor=north] at (\x,1.02) {\iv};}
  \foreach \n/\x/\s in {4/2.5/14, 5/6.5/10, 6/10.5/13, 7/14.5/12} {
    \node[boxnode, primaryblue, minimum size=8mm, font=\scriptsize] (t\n) at (\x,2.8) {\(t_{\n}{=}\s\)};}
  \node[boxnode, primaryblue, minimum size=8mm, font=\scriptsize] (t2) at (4.5,4.2) {\(t_{2}{=}24\)};
  \node[boxnode, primaryblue, minimum size=8mm, font=\scriptsize] (t3) at (12.5,4.2) {\(t_{3}{=}25\)};
  \node[boxnode, primaryblue, minimum size=9mm, line width=1pt, fill=primaryblue!16, font=\scriptsize] (t1) at (8.5,5.6) {\(t_{1}{=}49\)};
  % ---- edges ----
  \draw[flow] (t1) -- (t2); \draw[flow] (t1) -- (t3);
  \draw[flow] (t2) -- (t4); \draw[flow] (t2) -- (t5);
  \draw[flow] (t3) -- (t6); \draw[flow] (t3) -- (t7);
  \draw[flow] (t4) -- (t8); \draw[flow] (t4) -- (t9);
  \draw[flow] (t5) -- (t10); \draw[flow] (t5) -- (t11);
  \draw[flow] (t6) -- (t12); \draw[flow] (t6) -- (t13);
  \draw[flow] (t7) -- (t14); \draw[flow] (t7) -- (t15);
  \draw[flow] (t8) -- (1,-0.3); \draw[flow] (t8) -- (2,-0.3);
  \draw[flow] (t9) -- (3,-0.3); \draw[flow] (t9) -- (4,-0.3);
  \draw[flow] (t10) -- (5,-0.3); \draw[flow] (t10) -- (6,-0.3);
  \draw[flow] (t11) -- (7,-0.3); \draw[flow] (t11) -- (8,-0.3);
  \draw[flow] (t12) -- (9,-0.3); \draw[flow] (t12) -- (10,-0.3);
  \draw[flow] (t13) -- (11,-0.3); \draw[flow] (t13) -- (12,-0.3);
  \draw[flow] (t14) -- (13,-0.3); \draw[flow] (t14) -- (14,-0.3);
  \draw[flow] (t15) -- (15,-0.3); \draw[flow] (t15) -- (16,-0.3);
  \node[font=\scriptsize, text=darkgray, anchor=east] at (-0.2,0.1) {leaves \(A\)};
  \node[font=\scriptsize, text=darkgray, anchor=east] at (-0.2,5.6) {root};
\end{tikzpicture}
\caption{The sum segment tree over the running array \(A\) (bottom row: the
leaves, labelled by index \(1..16\)).  Every node \(t_i\) owns the interval
under it and stores the \emph{sum} of that interval; the green subscripts show
the eight size-2 nodes' intervals, and each larger interval is the union of its
children's.  Fold-up checks reused everywhere:
\(t_8{=}A_1{+}A_2{=}5\), \(t_4{=}t_8{+}t_9{=}5{+}9{=}14\),
\(t_2{=}t_4{+}t_5{=}14{+}10{=}24\) (the sum of slots \(1..8\), CSES~1648's first
query), and \(t_1{=}t_2{+}t_3{=}24{+}25{=}49=\sum A\).  Swap the fold for
\(\min\) or \(\max\) and the very same picture gives the tables below.  The
leaves \(t_{16..31}\) equal \(A\), so this array cross-checks with
Ch.18's~\figref{c14-array} at the leaf layer (the internal nodes differ:
interval aggregates, not lowbit blocks).}
\label{fig:c13-tree}
\end{figure}

The three folds over the internal nodes \(t_1..t_{15}\) (leaves \(t_{16..31}\)
are just \(A\)); every section re-uses this table and never recomputes it:

\begin{center}
\scriptsize
\renewcommand{\arraystretch}{1.1}
\setlength{\tabcolsep}{3pt}
\begin{tabular}{@{}l|cccc|cccc|ccccccc@{}}
\toprule
\textbf{node} & \(t_1\)&\(t_2\)&\(t_3\)&\(t_4\) & \(t_5\)&\(t_6\)&\(t_7\)&\(t_8\) & \(t_9\)&\(t_{10}\)&\(t_{11}\)&\(t_{12}\)&\(t_{13}\)&\(t_{14}\)&\(t_{15}\)\\
\midrule
\textbf{interval} & 1--16&1--8&9--16&1--4 & 5--8&9--12&13--16&1--2 & 3--4&5--6&7--8&9--10&11--12&13--14&15--16\\
\(\sum\) & 49&24&25&14 & 10&13&12&5 & 9&2&8&7&6&7&5\\
\(\min\) & 0&1&0&2 & 1&0&1&2 & 4&1&3&0&2&1&2\\
\(\max\) & 7&5&7&5 & 5&7&6&3 & 5&1&5&7&4&6&3\\
\bottomrule
\end{tabular}
\end{center}

The second running object is the tree \(T\) Chapters~12--17 have been sharing
-- nodes \(1..10\), root~1, edges
\[
  1\!-\!2,\; 1\!-\!3,\; 2\!-\!4,\; 2\!-\!5,\; 3\!-\!6,\;
  3\!-\!7,\; 5\!-\!8,\; 7\!-\!9,\; 9\!-\!10 ,
\]
with the Ch.16 preorder \(\tout\) (\emph{inclusive}):
\[
  \tin=(1,2,6,3,4,7,8,5,9,10),\qquad
  \tout=(10,5,10,3,5,7,10,5,10,10).
\]
So \(\subtree(2)=[2,5]\) holds \(\{2,4,5,8\}\) and \(\subtree(7)=[8,10]\)
holds \(\{7,9,10\}\) -- the intervals \secref{c13-static}'s Euler pointer hands
to a segment tree (Ch.~5 proved ``a subtree is one interval'',
\thmref{c24-range}; this chapter only supplies the range engine).

Five angles, five hand-computed overlays -- each lives in the section named on
its row and never changes:

\begin{center}
\small
\renewcommand{\arraystretch}{1.25}
\begin{tabular}{@{}p{0.22\linewidth} p{0.52\linewidth} l@{}}
\toprule
\textbf{Angle} & \textbf{What comes out} & \textbf{Traced in} \\
\midrule
fold leaves up & \(\min\)-tree root \(t_1{=}0\), \(\max\)-tree root \(t_1{=}7\), \(\sum\)-tree \(t_1{=}49\) -- one array, three aggregates, same shape & \secref{c13-core} \\
point-set cascade & set \(A_3{:}=9\) (was 4): only the nodes owning slot~3 re-fold -- \(t_9{=}14,\;t_4{=}19,\;t_2{=}29,\;t_1{=}54\); \(\min\) of \([1,4]\) still \(2\) (at \(A_2\)) & \secref{c13-core} \\
tuple merge & CSES~1190's \((1,2,-3,5,-1)\): best subarray \(=9\) for \([2,5]\); after set \(A_3{=}1\), \(=13\) & \secref{c13-combine} \\
tree descend & \(\max\)-tree over \(A\): first position with value \(\ge 7\) is slot~9 (\(A_9{=}7\)); found by walking down, not scanning & \secref{c13-search} \\
iterative \([l,r)\) & CF~514D's columns: fold the \(\max\) of a sliding window in the half-open two-accumulator walk; \(m\le5\) independent trees & \secref{c13-iter} \\
\bottomrule
\end{tabular}
\end{center}

% ---------------- Which aggregate, cost chart ----------------
\subsection*{Which Merge? (the aggregate chart)}

Every segment tree is the \emph{same} skeleton; the only thing that changes is
what one node stores and how two children merge.  That choice is forced by the
query shape -- and by whether the operation is even something a BIT could do.

\begin{center}
\small
\renewcommand{\arraystretch}{1.25}
\begin{tabular}{@{}p{0.15\linewidth} p{0.18\linewidth} p{0.21\linewidth} p{0.34\linewidth}@{}}
\toprule
\textbf{Query} & \textbf{Node stores} & \textbf{Identity} & \textbf{Why a segment tree (not a BIT)} \\
\midrule
range sum & one sum \(\bigO{1}\) merge & \(0\) & BIT does this too; segtree chosen for uniformity or non-sum siblings \\
range min/max & one value & \(+\infty / -\infty\) & \textbf{BIT cannot} do min/max under arbitrary \emph{writes} -- only append/decrement-only monotone cases \\
range gcd & one gcd & \(0\) & not invertible: no prefix trick exists \\
count in range & a count (freq.) & \(0\) & BIT can, but the tree descends for ``k-th'' \(\log n\) \\
max subarray & a 4-tuple (sum, pre, suf, best) & all \(-\infty\) & no BIT analogue: the merge must remember four things \\
first \(\ge x\) / k-th & any prefix-monotone aggregate & -- & \textbf{descend} the tree (19.3); the BIT cousin lifts from the other end \\
\bottomrule
\end{tabular}
\end{center}

\begin{insightbox}[One skeleton, and the whole game is ``what does a node remember'']
Rows 1--3 of the table are a single number per node; row~5 is a
\emph{four-field tuple}; row~6 is a scalar but read by walking \emph{down}
instead of folding \emph{up}.  When a segment-tree problem feels hard, you are
choosing a row -- deciding what each node must retain so that the merge stays
\(O(1)\) and associative -- not inventing a new data structure.  If the answer
is ``one scalar and the operation is a sum/count,'' ask first whether
Chapter~18's smaller BIT already fits.
\end{insightbox}

% ---------------- Taxonomy chart ----------------
\begin{figure}[H]
\centering
\begin{tikzpicture}[diagram, scale=0.95, transform shape,
    mnode/.style={circle, draw=#1, line width=0.9pt, fill=#1!8, minimum size=7mm,
                  font=\scriptsize\bfseries, text=#1!60!black}]
  \node[mnode=primaryblue] (a) at (0,0)     {interval};
  \node[annot, font=\scriptsize, anchor=west] at (1.25,0) {node = a range [l,r]; split at mid; leaf = one slot (19.1)};
  \node[mnode=successgreen] (b) at (0,-1)  {fold};
  \node[annot, font=\scriptsize, anchor=west] at (1.25,-1) {build = fold children up; update = retrace one root path; query = cover [ql,qr] by O(log n) nodes (19.1)};
  \node[mnode=difforange]   (c) at (0,-2)  {tuple};
  \node[annot, font=\scriptsize, anchor=west] at (1.25,-2) {node keeps (sum,pref,suff,best): max-subarray and any associative combine a scalar cannot hold (19.2)};
  \node[mnode=warningred]   (d) at (0,-3)  {descend};
  \node[annot, font=\scriptsize, anchor=west] at (1.25,-3) {walk root-to-leaf choosing a child: first >= x, k-th by counts -- binary search on the aggregate (19.3)};
  \node[mnode=softpurple]   (e) at (0,-4)  {2n form};
  \node[annot, font=\scriptsize, anchor=west] at (1.25,-4) {bottom-up iterative tree; half-open [l,r) query with two accumulators; minimal constant (19.4)};
  \node[mnode=deepblue]     (f) at (0,-5)  {boundary};
  \node[annot, font=\scriptsize, anchor=west] at (1.25,-5) {static RMQ -> Ch.21 sparse table; range-update -> Ch.20 lazy; path -> Ch.22 HLD (19.5)};
\end{tikzpicture}
\caption{The chapter's six-state taxonomy, read top to bottom as an
escalation.  Rows 1--2 are the machine (intervals, the three recursions);
row~3 is a change of \emph{what a node stores} (a tuple, not a scalar); row~4
runs the tree \emph{downward} (search instead of fold); row~5 is a change of
\emph{layout} (iterative \(2n\), no recursion); row~6 is a change of
\emph{boundary} -- knowing when this is a lazy (Ch.20), persistent (Ch.21), or
decomposition (Ch.22) problem after all.  Every solved problem sits on exactly
one row: CSES~1649 on rows~1--2, the CSES~1190 monoid on row~3, CF~1179C on
row~4, CF~514D on row~5, and the Summit spanning rows~2 and~4 (fold + descend).}
\label{fig:c13-taxonomy}
\end{figure}

% ---------------- Who owns what ----------------
\subsection*{Who Owns What (so nothing is learned twice)}

\begin{center}
\small
\renewcommand{\arraystretch}{1.25}
\begin{tabular}{@{}p{0.30\linewidth} p{0.30\linewidth} p{0.34\linewidth}@{}}
\toprule
\textbf{Idea} & \textbf{Owner} & \textbf{This chapter's share} \\
\midrule
lowbit blocks, prefix sums, difference BIT, k-th by lifting & Ch.~18 (\chref{c14}) & the invertible-op rival; 19.3 contrasts descend vs lift, 19.4 discharges the boundary table -- no BIT re-teach \\
Euler tour, \(\tin/\tout\), ``subtree \(=\) interval'' & Ch.~5 (\chref{c24}) & cited by label; 19.5 supplies the range \emph{engine} on the flat array \\
range \emph{update} + range query, lazy tags & Ch.~20 (\chref{c15}) & boundary is prose only; this chapter is point-update / range-query, lazy forward-cited \\
merge-sort tree, persistent tree, sparse-table RMQ & Ch.~21 (\chref{c19}) & static-RMQ ``there is a faster structure'' note is forward-cited; neither engine is built here \\
heavy-light path queries & Ch.~22 (\chref{c16}) & a segtree \emph{on} heavy paths lives there; mentioned by prose \\
binary lifting, LCA preorder & Ch.~11 (\chref{c11}) & 19.3 back-cites lifting as the ``descend by bits'' cousin; 19.5 uses Ch.16's \(\tin\) \\
\bottomrule
\end{tabular}
\end{center}

% ---------------- Recognition guide ----------------
\begin{recognitionbox}[Pattern card -- ``a range aggregate under point changes, where the aggregate is not a sum'']
\reckey{``set position \(k\) to \(u\) / add to position \(k\)'', interleaved
with ``\(\min\) / \(\max\) / gcd of \(x_a..x_b\)''; ``maximum subarray sum in
\([l,r]\) after point edits''; ``the first position whose value (or running
prefix) reaches \(x\)''; ``the \(k\)-th smallest in a multiset with
insert/delete''; any input whose brute force re-folds a range each query --
and, decisively, any \emph{non-invertible} fold (min/max/gcd) where a BIT has
no subtraction to lean on.}
\recuse{
  \item point mutations + a non-invertible range fold (min/max/gcd) -- the
        recursive tree, \secref{c13-core} (CSES~1649);
  \item the aggregate needs more than one scalar to merge (max subarray,
        longest run, mode-like) -- tuple nodes, \secref{c13-combine} (CSES~1190);
  \item the question is \emph{where} a threshold is crossed (first / k-th) --
        descend the tree, \secref{c13-search} (CF~1179C);
  \item you want the smallest constant / shortest code on a fixed-size array
        -- iterative \(2n\) form, \secref{c13-iter} (CF~514D).
}
\recnotuse{
  \item the operation is a plain sum/count under point updates -- Ch.~18's BIT
        (\chref{c14}) is smaller and faster; reach here only for uniformity or
        a non-invertible sibling;
  \item updates hit a whole \emph{range} (add to \([l,r]\), assign \([l,r]\))
        and queries ask ranges back -- that is lazy propagation, Ch.~20
        (\chref{c15}); this chapter stays point-update;
  \item the array never changes and you only ask \(\min/\max\) RMQs -- a
        sparse table (Ch.~21, \chref{c19}) answers in \(O(1)\); build a
        segment tree only if writes are possible;
  \item the ranges are tree \emph{paths}, not subtrees or array slices --
        heavy-light, Ch.~22 (\chref{c16}).
}
One-line test: \emph{is the question ``fold a range with an associative,
possibly non-invertible operation'' over a \emph{mutable} array, or ``find
where a prefix-monotone aggregate first crosses a threshold''?  Then a segment
tree is the right engine -- and this chapter is where its every variant is
chosen.}
\end{recognitionbox}

% ---------------- Constraint -> tool ----------------
\subsection*{Reading the Constraints (which form the limits want)}

\begin{center}
\small
\renewcommand{\arraystretch}{1.25}
\begin{tabular}{@{}p{0.32\linewidth} p{0.40\linewidth} p{0.22\linewidth}@{}}
\toprule
\textbf{Constraint signature} & \textbf{What it rules in/out} & \textbf{Tool} \\
\midrule
\(n,q\le2\cdot10^5\), point set + range min/max & prefix/BIT dead (not invertible); \(4n\) segtree is the workhorse & \secref{c13-core} \\
values up to \(10^9\), sum/gcd may overflow & node value is \code{long long}; \(\gcd\) of sums stays int but check & \secref{c13-core}, \secref{c13-combine} \\
range asks the \emph{best subarray}, with edits & scalar fold impossible -- tuple node \((\sum,\text{pre},\text{suf},\text{best})\) & \secref{c13-combine} (worked CSES 1190) \\
``first index where prefix \(\ge\mathit{target}\)'', online & binary search on prefix is \(O(\log^2 n)\); a tree descend is \(O(\log n)\) & \secref{c13-search} \\
fixed \(n\), tight time / code-golf, half-open ranges & iterative \(2n\) form: no recursion, best cache behaviour & \secref{c13-iter} (CF 514D) \\
array fully static, pure \(\min/\max\) RMQ, many queries & \(O(1)\) sparse table beats \(O(\log n)\) tree -- stop, go to Ch.~21 & \secref{c13-static} boundary \\
\bottomrule
\end{tabular}
\end{center}

% ---------------- Notation ----------------
\subsection*{Notation Used Only Here}

A node is named by its heap index \(t_v\) (root \(t_1\)), and \emph{owns} a
closed interval \([nl,nr]\) of array positions; \(\mid\) denotes its
\term{midpoint} \(mid=\lfloor(nl{+}nr)/2\rfloor\), so the left child owns
\([nl,mid]\) and the right owns \([mid{+}1,nr]\).  The \term{merge} (fold)
\(\mathcal{M}\) combines two children into their parent, and the
\term{identity} \(e\) is the neutral element for \(\mathcal{M}\) (an out-of-range
node contributes \(e\): \(0\) for sum, \(+\infty\) for min, \(-\infty\) for max,
\(0\) for gcd).  A \term{range query} \(\code{query}(q_l,q_r)\) returns
\(\mathcal{M}\) over the covering nodes and \emph{only} touches \(O(\log n)\)
of them; a \term{point update} \(\code{update}(p,u)\) sets \(A_p=u\) and
re-folds the single root path through \(p\).  The \term{iterative \(2n\) form}
(\secref{c13-iter}) stores leaves at \(t[n..2n)\) and parents at
\(t[i]=\mathcal{M}(t[2i],t[2i{+}1])\), answering a \emph{half-open}
\([l,r)\) query with a left and a right accumulator (order matters when
\(\mathcal{M}\) is not commutative).  To \term{descend} means to walk from the
root picking one child per level using a stored aggregate as the guide
(\secref{c13-search}); it is Chapter~1's binary search executed on the tree.
The Euler symbols \(\tin/\tout/\subtree(\cdot)\) are Ch.~5's, used in
\secref{c13-static} by reference only.  Every section runs on \(A\) and \(T\)
above: if a table here and a trace there disagree, this opener is authoritative
-- fix the section, never the other way.
